% ============================================================
%  GUÍA 4: ECUACIONES DE PRIMER GRADO EN ℤ
%  Curso Introductorio de Matemáticas — 1er Año
%  Autor: Ing. Gabriel Astudillo
%  Sesiones cubiertas: Semana 4 (clases 13 a 16)
% ============================================================

\documentclass[12pt, letterpaper]{article}

% ── Paquetes ─────────────────────────────────────────────────

\usepackage{mathwizards-guia}
\mwguia{Guía 4 — Ecuaciones de Primer Grado en $\Z$}{Ing. Gabriel Astudillo $\cdot$ Curso 1er Año}

% ── Comandos propios de esta guía ───────────────────
\newcommand{\Z}{\mathbb{Z}}

\begin{document}
% ============================================================

% ── Portada / Título ─────────────────────────────────────────
\mwportada{Guía de Estudio 4}{Ecuaciones de Primer Grado en $\Z$}{La Balanza Mágica: Igualdad, Lenguaje Algebraico y Resolución}

\vspace{4pt}
\begin{cuadroinfo}
  \small
  \textbf{Presentación:} Imagina una \textbf{balanza} en equilibrio: si agregas una piedra a un lado y no agregas nada en el otro, se desbalancea. Para mantenerla equilibrada, todo lo que hagas en un lado \emph{debes} hacerlo en el otro. Así funcionan las \textbf{ecuaciones}: son misterios donde una letra (la \textbf{incógnita}) esconde un número, y tú eres el detective que debe descubrirlo. En esta guía aprenderás a escribir problemas en «idioma matemático» y a resolver ecuaciones paso a paso, comprobando siempre tu respuesta.
  \\[4pt]
  \textbf{Nivel de Dificultad:}\quad
  \facil{} Básico / Directo \quad
  \medio{} Intermedio \quad
  \dificil{} Desafiante
\end{cuadroinfo}

\vspace{8pt}

% ============================================================
\section{La Ecuación y el Lenguaje Algebraico}
% ============================================================

\subsection{Igualdad, ecuación y sus partes}

Una \textbf{igualdad} es una frase matemática que dice que dos cosas valen lo mismo: $3 + 2 = 5$. Cuando en una igualdad hay una letra escondiendo un número, tenemos una \textbf{ecuación}.

\begin{cuadroteoria}
  \textbf{Definición:} una \textbf{ecuación} es una igualdad en la que aparece una o más \textbf{incógnitas} (letras) cuyo valor queremos descubrir.
  \[
    \underbrace{2x + 3}_{\text{1er miembro}} = \underbrace{11}_{\text{2do miembro}}
  \]
  \begin{itemize}
    \item \textbf{Miembros:} cada lado de la igualdad (izquierdo y derecho).
    \item \textbf{Incógnita:} la letra desconocida ($x$ aquí).
    \item \textbf{Término:} cada parte que se suma o se resta (aquí $2x$ y $3$).
    \item \textbf{Coeficiente:} el número que multiplica a la incógnita (aquí el $2$ de $2x$).
  \end{itemize}
\end{cuadroteoria}

\newpage

\subsection{Escribir en lenguaje algebraico}

Traducir frases al «idioma matemático» es el primer superpoder del algebrista:

\begin{cuadrosolucion}
  \textbf{Traducciones básicas} (con $x$ = «un número»):
  \begin{center}
    \begin{tabular}{ll}
      \textbf{Frase}              & \textbf{Lenguaje algebraico} \\ \hline
      un número aumentado en $5$  & $x + 5$                      \\
      un número disminuido en $3$ & $x - 3$                      \\
      el doble de un número       & $2x$                         \\
      el triple de un número      & $3x$                         \\
      la mitad de un número       & $\dfrac{x}{2}$               \\
      el cuadrado de un número    & $x^2$                        \\
    \end{tabular}
  \end{center}
\end{cuadrosolucion}

\begin{cuadroadvertencia}
  \textbf{Cuidado con el orden:} «un número menos $7$» es $x - 7$, pero «$7$ menos un número» es $7 - x$. \textquestiondown No son lo mismo!
\end{cuadroadvertencia}

\noindent\textbf{Ejemplo resuelto:} Escribe en lenguaje algebraico: «la edad de Ana es el doble de la edad de Luis».

Si la edad de Luis es $L$, entonces la de Ana es $2L$. Podemos escribir la igualdad $A = 2L$ (o directamente $2L$ si solo nos piden la expresión).

\subsection{Ejercicios: La Ecuación y el Lenguaje Algebraico}

\begin{enumerate}[leftmargin=*]
  \item \facil{} Identifica el primer miembro, el segundo miembro, la incógnita, todos los términos y el coeficiente de $x$ en la ecuación $-5x + 12 = -18$.

  \item \facil{} Escribe en lenguaje algebraico: «El triple de un número disminuido en $14$».

  \item \facil{} Escribe en lenguaje algebraico: «La cuarta parte de un número aumentada en $9$».

  \item \facil{} Escribe en lenguaje algebraico: «El quíntuple de un número menos $11$» y también «$11$ menos el quíntuple de un número». \textquestiondown Expresan lo mismo?

  \item \facil{} Escribe en lenguaje algebraico: «La suma de un número, su doble y su triple».

  \item \medio{} Escribe en lenguaje algebraico: «La mitad de la suma entre un número y $10$».

  \item \medio{} Traduce al lenguaje algebraico: «El cuádruple de la diferencia entre un número y $7$».

  \item \medio{} Escribe una expresión algebraica para: «La suma de tres números enteros consecutivos» (utiliza $n$ para el primer número).

  \item \medio{} Escribe una expresión algebraica para: «La suma de tres números enteros pares consecutivos» (utiliza $2n$ para el primer número par).

  \item \medio{} Traduce a una ecuación: «El triple de un número aumentado en $8$ equivale a su quíntuple disminuido en $12$».

  \item \medio{} Traduce a una ecuación: «Si al triple de la edad de Camila le restamos $5$ años, obtenemos $40$ años».

  \item \medio{} Escribe la expresión para el perímetro de un rectángulo cuyo largo mide $4$ cm más que el triple de su ancho $x$.

  \item \medio{} Traduce a una ecuación: «La tercera parte de un número disminuida en $6$ es igual a $-10$».

  \item \medio{} Traduce al lenguaje algebraico: «Pensé un número, lo multipliqué por $-3$, le sumé $15$ y obtuve $-6$».

  \item \medio{} Escribe en lenguaje algebraico: «El doble de un número disminuido en su tercera parte».

  \item \dificil{} Escribe la expresión para «la suma de los cuadrados de dos números $a$ y $b$» y compárala con «el cuadrado de la suma de $a$ y $b$». \textquestiondown Son equivalentes?

  \item \dificil{} Traduce con dos ecuaciones: «En una granja hay patos ($p$) y vacas ($v$). En total se cuentan $45$ cabezas y $130$ patas».

  \item \dificil{} Traduce con dos ecuaciones: «La edad de Don Roberto ($R$) supera en $25$ años a la de su hijo Andrés ($A$). Dentro de $10$ años, la edad del padre será el doble de la del hijo».

  \item \dificil{} Escribe en lenguaje algebraico y simplifica: «El cuádruple de un número, más el triple de su consecutivo, menos $8$».

  \item \dificil{} \textbf{:} Traduce a una ecuación con una sola incógnita: «Un padre tiene el triple de la edad de su hijo. Si la suma de sus edades más $5$ años es igual a $53$ años, \textquestiondown cuál es la ecuación que permite hallar la edad del hijo $h$?».
\end{enumerate}

%\newpage

% ============================================================
\section{Pasos para Resolver y Comprobar}
% ============================================================

\subsection{La balanza y las operaciones inversas}

Una ecuación es como una balanza en equilibrio. Para descubrir $x$:

\begin{cuadroteoria}
  \textbf{Regla de oro:} todo lo que hagas en un miembro, \textbf{hazlo también en el otro}. Así la balanza nunca se desequilibra.
\end{cuadroteoria}

\noindent Para «deshacer» lo que le pasa a la incógnita usamos la \textbf{operación inversa}:
\[
  \text{suma} \leftrightarrow \text{resta} \qquad\qquad \text{multiplicación} \leftrightarrow \text{división}
\]

\noindent\textbf{Ejemplo resuelto 1:} Resuelve $x + 5 = 12$.

A $x$ le sumaron $5$. Para deshacerlo, restamos $5$ en \emph{ambos} lados:
\[
  x + 5 - 5 = 12 - 5 \qquad \Rightarrow \qquad x = 7
\]

\noindent\textbf{Ejemplo resuelto 2:} Resuelve $3x = 21$.

A $x$ lo multiplicaron por $3$. Dividimos entre $3$ en ambos lados:
\[
  \frac{3x}{3} = \frac{21}{3} \qquad \Rightarrow \qquad x = 7
\]

\begin{cuadrosolucion}
  \textbf{La comprobación (¡nunca la saltes!):} sustituye el valor encontrado en la ecuación original y verifica que la igualdad se cumple.
  Para $x = 7$ en $x + 5 = 12$: $7 + 5 = 12$. \textquestiondown Se cumple! \checkmark
\end{cuadrosolucion}

\begin{cuadroadvertencia}
  \textbf{\textquestiondown La solución pertenece a $\Z$?} Como trabajamos con números enteros, revisa siempre si tu solución es un entero. Si te da $-3$ o $0$ o $12$, está en $\Z$; si te da algo con coma (como $2{,}5$), no pertenece a $\Z$ y hay que revisar el planteamiento.
\end{cuadroadvertencia}

\noindent\textbf{Ejemplo resuelto 3:} Resuelve $2x - 7 = 11$ y comprueba.

Sumamos $7$ en ambos lados: $2x = 18$. Dividimos entre $2$: $x = 9$. Comprobación: $2(9) - 7 = 18 - 7 = 11$. \checkmark

\subsection{Ejercicios: La Ecuación y el Lenguaje Algebraico}

\begin{enumerate}[leftmargin=*]
  \item \facil{} Identifica el primer miembro, el segundo miembro, la incógnita, todos los términos y el coeficiente de $x$ en la ecuación $-5x + 12 = -18$.

  \item \facil{} Escribe en lenguaje algebraico: «El triple de un número disminuido en $14$».

  \item \facil{} Escribe en lenguaje algebraico: «El cuádruple de un número aumentado en $9$».

  \item \facil{} Escribe en lenguaje algebraico: «El quíntuple de un número menos $11$» y también «$11$ menos el quíntuple de un número». \textquestiondown Expresan lo mismo?

  \item \facil{} Escribe en lenguaje algebraico: «La suma de un número, su doble y su triple».

  \item \medio{} Escribe en lenguaje algebraico: «El triple de la suma entre un número y $10$».

  \item \medio{} Traduce al lenguaje algebraico: «El cuádruple de la diferencia entre un número y $7$».

  \item \medio{} Escribe una expresión algebraica para: «La suma de tres números enteros consecutivos» (utiliza $n$ para el primer número).

  \item \medio{} Escribe una expresión algebraica para: «La suma de tres números enteros pares consecutivos» (utiliza $2n$ para el primer número par).

  \item \medio{} Traduce a una ecuación: «El triple de un número aumentado en $8$ equivale a su quíntuple disminuido en $12$».

  \item \medio{} Traduce a una ecuación: «Si al triple de la edad de Camila le restamos $5$ años, obtenemos $40$ años».

  \item \medio{} Escribe la expresión para el perímetro de un rectángulo cuyo largo mide $4$ cm más que el triple de su ancho $x$.

  \item \medio{} Traduce a una ecuación: «El quíntuple de un número disminuido en $6$ es igual a $-36$».

  \item \medio{} Traduce al lenguaje algebraico: «Pensé un número, lo multipliqué por $-3$, le sumé $15$ y obtuve $-6$».

  \item \medio{} Escribe en lenguaje algebraico: «El cuádruple de un número disminuido en su doble».

  \item \dificil{} Escribe la expresión para «la suma de los cuadrados de dos números $a$ y $b$» y compárala con «el cuadrado de la suma de $a$ y $b$». \textquestiondown Son equivalentes?

  \item \dificil{} Traduce con dos ecuaciones: «En una granja hay patos ($p$) y vacas ($v$). En total se cuentan $45$ cabezas y $130$ patas».

  \item \dificil{} Traduce con dos ecuaciones: «La edad de Don Roberto ($R$) supera en $25$ años a la de su hijo Andrés ($A$). Dentro de $10$ años, la edad del padre será el doble de la del hijo».

  \item \dificil{} Escribe en lenguaje algebraico y simplifica: «El cuádruple de un número, más el triple de su consecutivo, menos $8$».

  \item \dificil{} \textbf{:} Traduce a una ecuación con una sola incógnita: «Un padre tiene el triple de la edad de su hijo. Si la suma de sus edades más $5$ años es igual a $53$ años, \textquestiondown cuál es la ecuación que permite hallar la edad del hijo $h$?».
\end{enumerate}

%\newpage

% ============================================================
\section{Pasos para Resolver y Comprobar}
% ============================================================

\subsection{La balanza y las operaciones inversas}

Una ecuación es como una balanza en equilibrio. Para descubrir $x$:

\begin{cuadroteoria}
  \textbf{Regla de oro:} todo lo que hagas en un miembro, \textbf{hazlo también en el otro}. Así la balanza nunca se desequilibra.
\end{cuadroteoria}

\noindent Para «deshacer» lo que le pasa a la incógnita usamos la \textbf{operación inversa}:
\[
  \text{suma} \leftrightarrow \text{resta} \qquad\qquad \text{multiplicación} \leftrightarrow \text{división}
\]

\noindent\textbf{Ejemplo resuelto 1:} Resuelve $x + 5 = 12$.

A $x$ le sumaron $5$. Para deshacerlo, restamos $5$ en \emph{ambos} lados:
\[
  x + 5 - 5 = 12 - 5 \qquad \Rightarrow \qquad x = 7
\]

\noindent\textbf{Ejemplo resuelto 2:} Resuelve $3x = 21$.

A $x$ lo multiplicaron por $3$. Dividimos entre $3$ en ambos lados:
\[
  \frac{3x}{3} = \frac{21}{3} \qquad \Rightarrow \qquad x = 7
\]

\begin{cuadrosolucion}
  \textbf{La comprobación (¡nunca la saltes!):} sustituye el valor encontrado en la ecuación original y verifica que la igualdad se cumple.
  Para $x = 7$ en $x + 5 = 12$: $7 + 5 = 12$. \textquestiondown Se cumple! \checkmark
\end{cuadrosolucion}

\begin{cuadroadvertencia}
  \textbf{\textquestiondown La solución pertenece a $\Z$?} Como trabajamos con números enteros, revisa siempre si tu solución es un entero. Si te da $-3$ o $0$ o $12$, está en $\Z$; si te da algo con coma (como $2{,}5$), no pertenece a $\Z$ y hay que revisar el planteamiento.
\end{cuadroadvertencia}

\noindent\textbf{Ejemplo resuelto 3:} Resuelve $2x - 7 = 11$ y comprueba.

Sumamos $7$ en ambos lados: $2x = 18$. Dividimos entre $2$: $x = 9$. Comprobación: $2(9) - 7 = 18 - 7 = 11$. \checkmark

\subsection{Ejercicios: Pasos para Resolver y Comprobar}

\begin{enumerate}[leftmargin=*]
  \item \facil{} Resuelve y comprueba: $x - 14 = -25$

  \item \facil{} Resuelve y comprueba: $-6x = 42$

  \item \facil{} Resuelve: $-7x = -49$

  \item \facil{} Resuelve: $18 - x = 25$

  \item \medio{} Resuelve y comprueba: $3x - 17 = -5$

  \item \medio{} Resuelve y comprueba: $-5x + 8 = -27$

  \item \medio{} Resuelve: $-4x + 15 = -9$

  \item \medio{} Resuelve: $9 - 4x = -15$

  \item \medio{} Resuelve: $-2x + 11 = 3$

  \item \medio{} Resuelve: $4x - 18 = -50$

  \item \medio{} Resuelve: $7x + 15 = 2x - 10$

  \item \medio{} Resuelve: $-3x - 9 = 2x + 16$

  \item \medio{} Resuelve: $12 - 5x = 3x - 20$

  \item \medio{} Resuelve: $2x - 15 = -29$

  \item \medio{} Resuelve: $8 - 3x = -1$

  \item \dificil{} Resuelve y comprueba en $\Z$: $5x - 14 = 36$

  \item \dificil{} Resuelve: $-8x + 3 = -3x - 22$

  \item \dificil{} Resuelve: $-6x + 15 = -33$

  \item \dificil{} «Al multiplicar un número por $-4$ y restarle $10$, se obtiene $-42$.» \textquestiondown Cuál es el número? \textquestiondown Pertenece la solución a $\Z$?

  \item \dificil{} \textbf{:}
        \begin{enumerate}[label=\alph*)]
          \item Un estudiante resolvió $-2x + 6 = -10$ haciendo $-2x = -4 \Rightarrow x = 2$. \textquestiondown En qué paso se equivocó? \textquestiondown Cuál es la respuesta correcta?
          \item Resuelve y comprueba: $-3x + 14 = 35$.
        \end{enumerate}
\end{enumerate}

% ============================================================
\section{Ecuaciones Sin Signos de Agrupación}
% ============================================================

\subsection{El método completo}

Cuando la ecuación tiene varios términos, seguimos estos pasos:

\begin{cuadroteoria}
  \textbf{Pasos de resolución:}
  \begin{enumerate}
    \item \textbf{Simplifica} cada miembro (junta términos parecidos si puedes).
    \item \textbf{Junta las incógnitas} en un solo miembro (resta o suma el término con $x$ en ambos lados).
    \item \textbf{Junta los números} en el otro miembro.
    \item \textbf{Despeja:} divide por el coeficiente de $x$.
    \item \textbf{Comprueba} en la ecuación original y revisa si la solución está en $\Z$.
  \end{enumerate}
\end{cuadroteoria}

\noindent\textbf{Ejemplo resuelto:} Resuelve $5x - 7 = 3x + 9$.

Paso 2: restamos $3x$ en ambos lados: $5x - 7 - 3x = 9$, o sea $2x - 7 = 9$.
Paso 3: sumamos $7$: $2x = 16$.
Paso 4: dividimos entre $2$: $x = 8$.
Paso 5: comprobación: $5(8) - 7 = 40 - 7 = 33$ y $3(8) + 9 = 24 + 9 = 33$. \checkmark\ $8 \in \Z$.

\noindent\textbf{Ejemplo resuelto 2:} Resuelve $10 - x = 4$.

Restamos $10$: $-x = -6$. Multiplicamos por $-1$: $x = 6$. Comprobación: $10 - 6 = 4$. \checkmark

\subsection{Ejercicios: Ecuaciones Sin Signos de Agrupación}

\begin{enumerate}[leftmargin=*]
  \item \facil{} Resuelve: $4x - 15 = 17$

  \item \facil{} Resuelve: $-3x + 14 = -1$

  \item \facil{} Resuelve: $7x - 8 = 3x + 12$

  \item \facil{} Resuelve: $15 - 2x = 3$

  \item \facil{} Resuelve: $9x + 6 = 4x - 19$

  \item \medio{} Resuelve juntando términos semejantes primero: $5x - 9 + 2x = 3x + 15$

  \item \medio{} Resuelve: $8 - 5x = 2x - 13$

  \item \medio{} Resuelve: $-4x + 7 - x = 22 - 2x$

  \item \medio{} Resuelve: $9x - 12 - 4x = -2x + 23$

  \item \medio{} Resuelve: $10x - 3 = 14x + 21$

  \item \medio{} Resuelve: $8x - 15 = 3x + 20$

  \item \medio{} Resuelve: $7x - 18 = 3x + 22$

  \item \medio{} Resuelve: $6 - 3x + 8 = x - 14$

  \item \medio{} Resuelve: $-7x - 4 = -2x + 26$

  \item \medio{} Resuelve: $11x + 8 - 5x = 3x - 13$

  \item \dificil{} Resuelve: $9x - 14 = 4x + 21$

  \item \dificil{} Resuelve: $4x - 18 + 6x = 15 - 3x + 19$

  \item \dificil{} Resuelve: $11x - 19 = 4x + 23$

  \item \dificil{} «El quíntuple de un número, disminuida en $14$, es igual al triple del mismo número más $18$.» \textquestiondown Cuál es el número? \textquestiondown Pertenece a $\Z$?

  \item \dificil{} \textbf{:}
        \begin{enumerate}[label=\alph*)]
          \item Resuelve: $-5x + 18 - 2x = 4 - 9x + 26$ y comprueba tu solución.
          \item Determina si $x = -4$ es solución de $7 - 3x = 5x + 39$.
        \end{enumerate}
\end{enumerate}

%\newpage

% ============================================================
\section{Ecuaciones con Signos de Agrupación y Problemas de la Vida Real}
% ============================================================

\subsection{Primero quita los paréntesis}

Cuando aparecen paréntesis, el primer paso es \textbf{eliminarlos} usando la propiedad distributiva:

\begin{cuadroteoria}
  \textbf{Propiedad distributiva:}
  \[
    a(x + b) = ax + ab \qquad\qquad -(x + b) = -x - b
  \]
  \textquestiondown Cuidado con el signo menos delante del paréntesis! Cambia todos los signos de adentro.
\end{cuadroteoria}

\noindent\textbf{Ejemplo resuelto 1:} Resuelve $2(x + 3) = 14$.

Distribuimos: $2x + 6 = 14$. Restamos $6$: $2x = 8$. Dividimos: $x = 4$.
Comprobación: $2(4 + 3) = 2 \times 7 = 14$. \checkmark

\noindent\textbf{Ejemplo resuelto 2 (problema):} «Pensé un número, lo tripliqué, le sumé $7$ y obtuve $31$. \textquestiondown Cuál es el número?»

Sea $x$ el número. La ecuación es $3x + 7 = 31$. Restamos $7$: $3x = 24$. Dividimos: $x = 8$.
Comprobación: $3(8) + 7 = 24 + 7 = 31$. \checkmark\ El número es $8$.

\begin{cuadrosolucion}
  \textbf{Los 6 pasos del resolutor de problemas:}
  \begin{enumerate}
    \item Lee el problema con calma y Elige la incógnita (¿qué me preguntan?).
    \item Escribe la ecuación.
    \item Resuélvela.
    \item Comprueba la solución.
    \item Responde con una frase completa.
  \end{enumerate}
\end{cuadrosolucion}

\subsection{Ejercicios: Ecuaciones con Signos de Agrupación y Problemas de la Vida Real}

\begin{enumerate}[leftmargin=*]
  \item \facil{} Resuelve: $3(x - 4) = 15$

  \item \facil{} Resuelve: $-2(x + 5) = -16$

  \item \facil{} Resuelve: $4(2x - 1) = 28$

  \item \facil{} Resuelve: $5(x + 3) - 7 = 23$

  \item \medio{} Resuelve: $5(x - 2) = 2(x + 7)$

  \item \medio{} Resuelve: $-4(x - 3) + 2x = 3(x - 1) - 5$

  \item \medio{} Resuelve: $2(3x + 1) - (x - 4) = 21$

  \item \medio{} Resuelve: $6(x - 2) - 2(x - 5) = 14$

  \item \medio{} Resuelve: $3[x - 2(x - 3)] = 12$

  \item \dificil{} Resuelve: $2[4 - 3(x - 2)] = 4(x + 5)$

  \item \dificil{} Resuelve: $5(2 - x) - 3[2x - (x + 1)] = -19$

  \item \medio{} \textbf{Problema de números consecutivos:} La suma de tres números enteros consecutivos es igual a $-48$. \textquestiondown Cuáles son los tres números?

  \item \medio{} \textbf{Problema de compras y dinero:} Sofía compró un cuaderno que costó el triple que un lápiz. Si por ambos artículos pagó $\$4.800$, \textquestiondown cuánto costó cada producto?

  \item \medio{} \textbf{Problema de edades (presente):} La edad de Carlos supera en $7$ años al doble de la edad de su primo Mateo. Si la suma de sus edades es $37$ años, \textquestiondown qué edad tiene Mateo? \textquestiondown Y Carlos?

  \item \medio{} \textbf{Problema de geometría (perímetro):} El largo de un terreno rectangular mide $8$ metros más que su ancho. Si el perímetro del terreno es de $64$ metros, calcula las dimensiones del terreno (ancho y largo).

  \item \medio{} \textbf{Problema de reparto:} Entre tres amigos (Ana, Beto y Carla) juntaron $150$ estampillas. Beto reunió el doble que Ana, y Carla reunió $10$ estampillas más que Beto. \textquestiondown Cuántas estampillas reunió cada uno?

  \item \dificil{} \textbf{Problema de edades en el futuro:} El señor Fernández tiene $42$ años y su hija Valeria tiene $12$ años. \textquestiondown Dentro de cuántos años la edad del padre será el doble de la de su hija?

  \item \dificil{} \textbf{Problema de geometría (triángulo):} En un triángulo isósceles, cada uno de los lados iguales mide $5$ cm más que la base. Si el perímetro del triángulo es $43$ cm, \textquestiondown cuánto mide la base y cuánto cada lado igual?

  \item \dificil{} \textbf{Problema de alcancía y monedas:} Tomás tiene en su alcancía monedas de $\$100$ y de $\$500$. Si tiene en total $24$ monedas que suman $\$6.400$, \textquestiondown cuántas monedas de cada tipo tiene en la alcancía?

  \item \dificil{} \textbf{Problema de números pares consecutivos:} La suma de tres números enteros pares consecutivos es $-102$. Encuentra los tres números.

  \item \dificil{} \textbf{ 1 (Adivina el número):} «Pensé un número entero, lo multipliqué por $-3$, al resultado le sumé $18$, luego multipliqué todo por $2$ y finalmente le resté $10$, obteniendo $32$». \textquestiondown Cuál fue el número que pensé? (Escribe la ecuación completa con paréntesis y corchetes, resuélvela paso a paso y comprueba).

  \item \dificil{} \textbf{ 2 (Desafío de edades cruzadas):} Un padre tiene $45$ años y sus dos hijos tienen $8$ y $12$ años. \textquestiondown Dentro de cuántos años la edad del padre será igual a la suma de las edades de sus dos hijos?
\end{enumerate}

\newpage

% ============================================================
\ifmwclaves
  \section{Clave de Respuestas}
  % ============================================================

  \subsection*{Lenguaje Algebraico (Ejercicios 1--20)}

  \begin{enumerate}[leftmargin=*, label=\textbf{\arabic*.}]
    \item 1er miembro: $-5x + 12$, 2do miembro: $-18$. Incógnita: $x$. Términos: $-5x$, $12$, $-18$. Coeficiente de $x$: $-5$.
    \item $3x - 14$
    \item $4x + 9$
    \item $5x - 11$ y $11 - 5x$. \textquestiondown No expresan lo mismo! (El orden de la resta cambia el resultado).
    \item $x + 2x + 3x$ (o simplificado $6x$).
    \item $3(x + 10)$
    \item $4(x - 7)$
    \item $n + (n + 1) + (n + 2)$
    \item $2n + (2n + 2) + (2n + 4)$
    \item $3x + 8 = 5x - 12$
    \item $3c - 5 = 40$ (o bien $3x - 5 = 40$).
    \item Ancho $= x$, Largo $= 3x + 4$. Perímetro $= 2(x + 3x + 4) = 8x + 8$.
    \item $5x - 6 = -36$
    \item $-3x + 15 = -6$
    \item $4x - 2x$ (o simplificado $2x$).
    \item Suma de cuadrados: $a^2 + b^2$; Cuadrado de la suma: $(a + b)^2$. \textquestiondown No son equivalentes! Por ejemplo, si $a=1$ y $b=2$, $1^2+2^2=5$ pero $(1+2)^2=9$.
    \item $p + v = 45$ \quad y \quad $2p + 4v = 130$
    \item $R = A + 25$ \quad y \quad $R + 10 = 2(A + 10)$
    \item $4x + 3(x + 1) - 8 = 4x + 3x + 3 - 8 = 7x - 5$
    \item $h + 3h + 5 = 53$ \quad (o bien $4h + 5 = 53$).
  \end{enumerate}

  \subsection*{Pasos para Resolver (Ejercicios 21--40)}

  \begin{enumerate}[leftmargin=*, label=\textbf{\arabic*.}]
    \setcounter{enumi}{20}
    \item $x = -25 + 14 \Rightarrow x = -11$. \quad (Comprobación: $-11 - 14 = -25$ \checkmark)
    \item $x = \dfrac{42}{-6} \Rightarrow x = -7$. \quad (Comprobación: $-6(-7) = 42$ \checkmark)
    \item $-7x = -49 \Rightarrow x = 7$
    \item $-x = 25 - 18 = 7 \Rightarrow x = -7$
    \item $3x = -5 + 17 = 12 \Rightarrow x = 4$. \quad (Comprobación: $3(4) - 17 = -5$ \checkmark)
    \item $-5x = -27 - 8 = -35 \Rightarrow x = 7$. \quad (Comprobación: $-5(7) + 8 = -27$ \checkmark)
    \item $-4x = -9 - 15 = -24 \Rightarrow x = 6$
    \item $-4x = -15 - 9 = -24 \Rightarrow x = 6$
    \item $-2x = 3 - 11 = -8 \Rightarrow x = 4$
    \item $4x = -50 + 18 = -32 \Rightarrow x = -8$
    \item $7x - 2x = -10 - 15 \Rightarrow 5x = -25 \Rightarrow x = -5$
    \item $-3x - 2x = 16 + 9 \Rightarrow -5x = 25 \Rightarrow x = -5$
    \item $-5x - 3x = -20 - 12 \Rightarrow -8x = -32 \Rightarrow x = 4$
    \item $2x = -29 + 15 = -14 \Rightarrow x = -7$
    \item $-3x = -1 - 8 = -9 \Rightarrow x = 3$
    \item $5x = 36 + 14 = 50 \Rightarrow x = 10 \in \Z$. \quad (Comprobación: $5(10) - 14 = 36$ \checkmark)
    \item $-8x + 3x = -22 - 3 \Rightarrow -5x = -25 \Rightarrow x = 5$
    \item $-6x = -33 - 15 = -48 \Rightarrow x = 8$
    \item $-4x - 10 = -42 \Rightarrow -4x = -32 \Rightarrow x = 8 \in \Z$ \checkmark
    \item a) Al trasladar $+6$ al segundo miembro olvidó cambiar el signo (debía restar $6$): $-2x = -10 - 6 = -16 \Rightarrow x = 8$. \quad b) $-3x = 35 - 14 = 21 \Rightarrow x = -7$. (Comprobación: $-3(-7) + 14 = 21 + 14 = 35$ \checkmark)
  \end{enumerate}

  \subsection*{Sin Signos de Agrupación (Ejercicios 41--60)}

  \begin{enumerate}[leftmargin=*, label=\textbf{\arabic*.}]
    \setcounter{enumi}{40}
    \item $4x = 17 + 15 = 32 \Rightarrow x = 8$
    \item $-3x = -1 - 14 = -15 \Rightarrow x = 5$
    \item $7x - 3x = 12 + 8 \Rightarrow 4x = 20 \Rightarrow x = 5$
    \item $-2x = 3 - 15 = -12 \Rightarrow x = 6$
    \item $9x - 4x = -19 - 6 \Rightarrow 5x = -25 \Rightarrow x = -5$
    \item $7x - 9 = 3x + 15 \Rightarrow 4x = 24 \Rightarrow x = 6$
    \item $-5x - 2x = -13 - 8 \Rightarrow -7x = -21 \Rightarrow x = 3$
    \item $-5x + 7 = 22 - 2x \Rightarrow -3x = 15 \Rightarrow x = -5$
    \item $5x - 12 = -2x + 23 \Rightarrow 7x = 35 \Rightarrow x = 5$
    \item $10x - 14x = 21 + 3 \Rightarrow -4x = 24 \Rightarrow x = -6$
    \item $8x - 3x = 20 + 15 \Rightarrow 5x = 35 \Rightarrow x = 7$
    \item $7x - 3x = 22 + 18 \Rightarrow 4x = 40 \Rightarrow x = 10$
    \item $14 - 3x = x - 14 \Rightarrow -4x = -28 \Rightarrow x = 7$
    \item $-7x + 2x = 26 + 4 \Rightarrow -5x = 30 \Rightarrow x = -6$
    \item $6x + 8 = 3x - 13 \Rightarrow 3x = -21 \Rightarrow x = -7$
    \item $9x - 4x = 21 + 14 \Rightarrow 5x = 35 \Rightarrow x = 7$
    \item $10x - 18 = 34 - 3x \Rightarrow 13x = 52 \Rightarrow x = 4$
    \item $11x - 4x = 23 + 19 \Rightarrow 7x = 42 \Rightarrow x = 6$
    \item $5x - 14 = 3x + 18 \Rightarrow 2x = 32 \Rightarrow x = 16 \in \Z$ \checkmark
    \item a) $-7x + 18 = -9x + 30 \Rightarrow 2x = 12 \Rightarrow x = 6$. Comprobación: $-5(6)+18-2(6) = -24$ y $4-9(6)+26 = -24$ \checkmark \quad b) Sustituyendo $x = -4$: $7 - 3(-4) = 19$ y $5(-4) + 39 = 19$. ¡Sí es solución!
  \end{enumerate}

  \subsection*{Con Signos de Agrupación y Problemas de la Vida Real (Ejercicios 61--82)}

  \begin{enumerate}[leftmargin=*, label=\textbf{\arabic*.}]
    \setcounter{enumi}{60}
    \item $3x - 12 = 15 \Rightarrow 3x = 27 \Rightarrow x = 9$
    \item $-2x - 10 = -16 \Rightarrow -2x = -6 \Rightarrow x = 3$
    \item $8x - 4 = 28 \Rightarrow 8x = 32 \Rightarrow x = 4$
    \item $5x + 15 - 7 = 23 \Rightarrow 5x + 8 = 23 \Rightarrow 5x = 15 \Rightarrow x = 3$
    \item $5x - 10 = 2x + 14 \Rightarrow 3x = 24 \Rightarrow x = 8$
    \item $-4x + 12 + 2x = 3x - 3 - 5 \Rightarrow -2x + 12 = 3x - 8 \Rightarrow -5x = -20 \Rightarrow x = 4$
    \item $6x + 2 - x + 4 = 21 \Rightarrow 5x + 6 = 21 \Rightarrow 5x = 15 \Rightarrow x = 3$
    \item $6x - 12 - 2x + 10 = 14 \Rightarrow 4x - 2 = 14 \Rightarrow 4x = 16 \Rightarrow x = 4$
    \item $3[x - 2x + 6] = 3[-x + 6] = -3x + 18 = 12 \Rightarrow -3x = -6 \Rightarrow x = 2$
    \item $2[10 - 3x] = 20 - 6x = 4x + 20 \Rightarrow -10x = 0 \Rightarrow x = 0$
    \item $10 - 5x - 3[x - 1] = 10 - 5x - 3x + 3 = 13 - 8x = -19 \Rightarrow -8x = -32 \Rightarrow x = 4$
    \item Ecuación: $n + (n+1) + (n+2) = -48 \Rightarrow 3n + 3 = -48 \Rightarrow 3n = -51 \Rightarrow n = -17$. Los números son $-17$, $-16$ y $-15$.
    \item Lápiz $x$, cuaderno $3x$. Ecuación: $x + 3x = 4800 \Rightarrow 4x = 4800 \Rightarrow x = 1200$. Lápiz: $\$1.200$, Cuaderno: $\$3.600$.
    \item Mateo $m$, Carlos $2m + 7$. Ecuación: $m + (2m + 7) = 37 \Rightarrow 3m = 30 \Rightarrow m = 10$. Mateo tiene $10$ años y Carlos $27$ años.
    \item Ancho $x$, largo $x + 8$. Ecuación: $2(x + x + 8) = 64 \Rightarrow 4x + 16 = 64 \Rightarrow 4x = 48 \Rightarrow x = 12$. Ancho: $12$ m, Largo: $20$ m.
    \item Ana $a$, Beto $2a$, Carla $2a + 10$. Ecuación: $a + 2a + 2a + 10 = 150 \Rightarrow 5a = 140 \Rightarrow a = 28$. Ana: $28$, Beto: $56$, Carla: $66$ estampillas.
    \item Años $t$. Ecuación: $42 + t = 2(12 + t) \Rightarrow 42 + t = 24 + 2t \Rightarrow t = 18$ años.
    \item Base $b$, lados iguales $b + 5$. Ecuación: $b + 2(b + 5) = 43 \Rightarrow 3b + 10 = 43 \Rightarrow 3b = 33 \Rightarrow b = 11$ cm. Base: $11$ cm, cada lado igual: $16$ cm.
    \item Monedas de $\$500$: $x$, de $\$100$: $24 - x$. Ecuación: $500x + 100(24 - x) = 6400 \Rightarrow 400x + 2400 = 6400 \Rightarrow 400x = 4000 \Rightarrow x = 10$. Monedas de $\$500$: $10$, de $\$100$: $14$.
    \item Ecuación: $2n + (2n + 2) + (2n + 4) = -102 \Rightarrow 6n + 6 = -102 \Rightarrow 6n = -108 \Rightarrow n = -18$. Los números son $-36$, $-34$ y $-32$.
    \item Ecuación: $2(-3x + 18) - 10 = 32 \Rightarrow -6x + 36 - 10 = 32 \Rightarrow -6x + 26 = 32 \Rightarrow -6x = 6 \Rightarrow x = -1$. (Comprobación: $2(-3(-1) + 18) - 10 = 2(21) - 10 = 42 - 10 = 32$ \checkmark)
    \item Años $x$. Ecuación: $45 + x = (8 + x) + (12 + x) \Rightarrow 45 + x = 20 + 2x \Rightarrow x = 25$ años. (Comprobación: $45 + 25 = 70$ y $33 + 37 = 70$ \checkmark)
  \end{enumerate}

  \vspace{10pt}
  \begin{cuadrosolucion}
    \textbf{\textquestiondown Cómo te fue?} \textquestiondown Descubriste todos los números escondidos? Eres oficialmente un resolutor de misterios matemáticos. \textquestiondown Nos vemos en el \textbf{examen final} para demostrar todo lo que aprendiste!
  \end{cuadrosolucion}

\fi
\end{document}
